% FREUID Challenge 2026 — Technical Report
% Compile: pdflatex freuid_technical_report.tex && bibtex freuid_technical_report && pdflatex x2
% License: CC0 template (organizer-provided) — content below is our own.

\documentclass[11pt,a4paper]{article}

\usepackage[T1]{fontenc}
\usepackage[utf8]{inputenc}
\usepackage{lmodern}
\usepackage{microtype}
\usepackage{geometry}
\usepackage{graphicx}
\usepackage{booktabs}
\usepackage{amsmath}
\usepackage{hyperref}
\usepackage{xcolor}
\usepackage{enumitem}

\geometry{margin=2.5cm}
\hypersetup{
  colorlinks=true,
  linkcolor=blue!60!black,
  citecolor=blue!60!black,
  urlcolor=blue!60!black,
}

\newcommand{\teamname}{Nadhir Hasan}
\newcommand{\kaggleteam}{Nadhir Hasan}
\newcommand{\contactemail}{hasannadhir2003@gmail.com}
\newcommand{\reportdate}{\today}

\title{%
  \textbf{The FREUID Challenge 2026}\\[0.4em]
  \large Technical Report --- \teamname
}
\author{%
  Nadhir Hasan\\
  \small
  \texttt{\contactemail}
}
\date{Report date: \reportdate}

\begin{document}
\maketitle

\begin{abstract}
We present our solution to the FREUID Challenge 2026 identity-document fraud
detection task. Our final model is a DINOv2-L Vision Transformer fine-tuned with
low-rank adapters (LoRA, rank 16) and a per-patch multiple-instance-learning
classification head, trained at 448$\times$728 resolution. We train jointly on the
official FREUID training set (5 document types) and 8 countries from the
externally licensed IDNet dataset, using an annotation-driven synthetic fraud
generator (portrait swap, cross-card region swap, text-field editing, and
erase-and-retype) built on manually labeled per-type face and text-field bounding
boxes. Critically, we select the training checkpoint using the exact competition
metric (FREUID Score, combining AuDET and APCER@1\%BPCER) evaluated on two IDNet
countries held out entirely from training, rather than an in-domain FREUID
validation split, which we found during development to be an unreliable predictor
of leaderboard rank due to asset/persona duplication within the FREUID training
distribution. Our submitted model scores [0.xxxxx] on the public leaderboard while
maintaining substantially better cross-document-type generalization than a
FREUID-only-trained variant of the same architecture. Code, weights, and a
network-isolated Docker reproduction container are provided in the linked
repository.
\end{abstract}

\noindent\textbf{Competition:} The FREUID Challenge 2026 (IJCAI-ECAI 2026)\\
\textbf{Kaggle team:} \kaggleteam\\
\textbf{Code:} \url{https://github.com/nadhirhasan/FREUID-Challenge-2026} (commit \texttt{b9473882f66a8baefce83980cf19f4db23d11a33})\\

\section{Introduction}
\label{sec:intro}
Identity-document fraud detection must generalize across a threat model that spans
physical tampering, GenAI-driven digital manipulation, and print-and-capture
(``analog hole'') attacks, and across document types and issuing countries that are
not exhaustively represented in any single training set. The FREUID Challenge
formalizes this with a training set covering 5 document types and a scoring metric
(FREUID Score, Section~\ref{sec:results}) that combines a global ranking measure
(AuDET) with a strict low-false-alarm operating point (APCER@1\%BPCER), penalizing
solutions that overfit to easy global separability while failing at the
operating point that matters for production KYC systems~\cite{freuid2026}.

Our approach is built around a specific empirical finding from our own
development process: a classifier trained \emph{only} on the FREUID distribution
can reach extremely low FREUID Score on the public leaderboard while being
statistically indistinguishable from random (AUC $\approx 0.51$) on document types
it has never seen. We treat this as a direct proxy for the risk the private test
set is designed to expose (it is drawn from 7 document types, including types not
present in the public training data). We therefore deliberately trade some
public-leaderboard score for measured cross-type generalization, by mixing in an
external, permissively licensed ID-document dataset (IDNet) during training and
selecting checkpoints on a held-out slice of that external data rather than on
any FREUID-derived validation signal.

\section{Method}
\label{sec:method}

\subsection{Overview}
Each input image is loaded at full resolution, converted to RGB, letterboxed
(aspect-ratio-preserving resize + zero-padding) to a fixed $448\times728$ canvas,
and normalized with standard ImageNet statistics. The backbone produces a token
sequence; a per-patch head scores each spatial token independently and aggregates
via a combination of top-$k$ mean and attention-weighted pooling into a single
image-level fraud logit. A sigmoid converts this logit to the submitted score in
$[0,1]$; no further post-processing or calibration is applied (Section~\ref{sec:calibration}).

\subsection{Model architecture}
Backbone: \texttt{vit\_large\_patch14\_reg4\_dinov2} (DINOv2-L~\cite{dinov2} with
register tokens), loaded via \texttt{timm}, Apache-2.0 licensed pre-trained weights, 24
transformer blocks, embedding dimension 1024. We attach a from-scratch, minimal
LoRA implementation (rank 16, $\alpha=32$, dropout 0.05; no third-party
\texttt{peft} dependency) to the \texttt{qkv}, \texttt{proj}, \texttt{fc1}, and
\texttt{fc2} linear layers of every transformer block; the backbone itself is
frozen. On top of the token sequence we place a lightweight per-patch
multiple-instance-learning (MIL) head: a 2-layer MLP (hidden size 512, GELU,
dropout 0.3) scores every non-prefix token, and the image-level logit is the mean
of (i) the top-10\% most suspicious patch logits and (ii) an attention-weighted
sum over all patch logits, each contributing equally. This design was chosen to
push the model toward \emph{local}, document-type-agnostic forgery cues (blend
seams, splice boundaries, compression discontinuities) rather than global,
country-specific layout statistics, which we found improves transfer to unseen
document types. Trainable parameters: LoRA adapters + head only ($\approx$7M of
$\approx$304M total). Training/inference batch size: 6 (accumulated to an
effective batch of 24) / 16--32 (inference only, no gradient state).

\subsection{Training objective and procedure}
Loss: Focal BCE ($\alpha=0.25$, $\gamma=2.0$) on the single whole-image logit.
Optimizer: AdamW, two parameter groups (head lr $1\times10^{-3}$, LoRA lr
$2\times10^{-4}$, weight decay 0.05/0 respectively), cosine learning-rate
schedule with 5\% linear warmup. Mixed precision (fp16 autocast + gradient
scaler). Trained for 3 epochs over $\approx$111{,}000 images/epoch on a single
NVIDIA RTX A4500 (20GB); $\approx$12--14 hours per fold at this configuration.
We fine-tune the DINOv2 pre-trained weights via LoRA (frozen backbone) rather
than training any component from scratch, except the LoRA adapters and the MIL
head, which are randomly initialized.

\textbf{Synthetic fraud generation.} 25\% of genuine training images (from both
FREUID and IDNet sources) are converted to synthetic fraud examples each epoch,
using one of four attack types, dispatched via manually annotated per-document-type
face and text-field bounding boxes (Section~\ref{sec:external-data}):
cross-card region/portrait swap (donor of the same document type), text-field
edit (same-line content shifted and re-blended), erase-and-retype (ink removed via
inpainting, new text rendered with a system font), and a generic self-blend
fallback. Blending uses a mix of feathered alpha compositing and seamless
(Poisson) cloning to avoid a single, learnable seam signature. This attack suite,
compared against a simpler untargeted self-blend-only baseline under an otherwise
identical recipe, improved our public leaderboard score by
$\approx$43\% in isolated ablation (Section~\ref{sec:results}).

\subsection{Score calibration}
\label{sec:calibration}
No score calibration or thresholding is applied: the submitted label is
$\sigma(\text{logit})\in(0,1)$ directly from the model's forward pass, identical
at validation and test time. We rely on the FREUID Score's monotone invariance to
score rescaling (it depends only on ranking and the genuine-score quantile that
sets the 1\%-BPCER threshold), so no train/test-specific calibration step is
necessary or applied.

\section{Data}
\label{sec:data}

\subsection{FREUID training set}
\label{sec:external-data}
We use the official \texttt{train/} images and \texttt{train\_labels.csv}
metadata, stratified into 5 folds by (\texttt{type} $\times$ \texttt{label}) with
a fixed seed; the submitted model uses fold 0 (all 5 FREUID document types are
present in training --- we do not hold out an entire document type, since the
public/private test sets are drawn from the same underlying document-type
distribution as the training set, per the challenge description). We additionally
hand-annotated the portrait and text-field bounding boxes for every document
type (5 FREUID + 10 IDNet countries) to drive the synthetic attack generator
precisely rather than at random image locations.

\subsection{External data and pre-trained models}
\begin{itemize}[noitemsep]
  \item \textbf{DINOv2}~\cite{dinov2} pre-trained ViT-L/14 weights, via
    \texttt{timm} (Apache-2.0 license).
  \item \textbf{IDNet}~\cite{idnet2024} (real, non-FREUID identity documents
    across 10 European countries), CC BY 4.0 license, publicly available at
    \url{https://arxiv.org/abs/2409.10472}. 8 countries (Spain, Albania,
    Azerbaijan, Finland, Greece, Latvia, Russia, Serbia) are mixed into
    training, capped at 56{,}000 images (balanced genuine/fraud) to roughly
    match the per-fold FREUID training volume. 2 further countries (Estonia,
    Slovakia) are held out entirely --- never seen in training in any form ---
    and used exclusively for checkpoint selection (Section~\ref{sec:validation}).
\end{itemize}
Both artifacts are publicly available under licenses compatible with this
competition's external-data policy (``any publicly available pre-trained model
or dataset may be used, provided the license is compatible and the artifact is
fully cited''); no proprietary or non-freely-accessible data was used.
We deliberately avoided any model, checkpoint, or code released under a
non-OSI-approved or non-commercial license (e.g.\ TruFor, ConvNeXt V2, DINOv3)
anywhere in the training or inference pipeline for this submission.

\subsection{Validation strategy}
\label{sec:validation}
During development we found that in-domain FREUID validation (a stratified,
disjoint 20\% fold) is an \emph{unreliable} predictor of public-leaderboard rank:
in one controlled comparison, the fold with the numerically best in-domain
validation score (by six orders of magnitude) scored \emph{worse} on the public
leaderboard than a fold with a substantially worse in-domain score. We traced
this to heavy visual asset/persona reuse within the FREUID training distribution
(near-duplicate faces and field values recur across nominally disjoint
train/validation splits), which lets a model ``pass'' in-domain validation by
recognizing recycled assets rather than genuine forgery cues. We therefore select
the checkpoint to submit using the exact competition FREUID Score
(Section~\ref{sec:results}) computed on the IDNet held-out slice described above
(Section~\ref{sec:external-data}), which shares no assets with any training data
by construction. Private test labels were never used, seen, or referenced during
model selection.

\section{Inference}
\label{sec:inference}
At inference, each image is loaded, letterboxed to $448\times728$ (identical to
training preprocessing, no cropping/tiling), and passed through the model in
batches (no test-time augmentation or ensembling in the submitted configuration).
The Docker entrypoint (\texttt{docker/prepare\_submission.py}) discovers every
image directly under \texttt{/data/}, runs inference, and writes
\texttt{/submissions/submission.csv} with columns \texttt{id,label}. The
container requires no network access at runtime (\texttt{--network none}); all
model weights are baked into the image at build time. Inference auto-detects
CUDA and falls back to CPU execution if unavailable.

\section{Results}
\label{sec:results}

\begin{table}[h]
  \centering
  \caption{Summary of main results (FREUID Score; lower is better).}
  \label{tab:results}
  \begin{tabular}{lcc}
    \toprule
    Split / Leaderboard & FREUID $\downarrow$ & Notes \\
    \midrule
    IDNet held-out (selection metric)       & 0.1028 & Selected checkpoint (epoch [X]) \\
    Public leaderboard                       & [0.xxxxx] & Selected final submission \\
    Private (if known)                       & --     & Not available at report time \\
    \bottomrule
  \end{tabular}
\end{table}

\noindent\textbf{Ablation --- attack realism.} Holding the architecture, data
split, and all other hyperparameters fixed, replacing an untargeted self-blend-only
synthetic-fraud generator with the annotation-driven attack suite
(Section~\ref{sec:method}) improved the public leaderboard score from 0.00169 to
0.00096 (a FREUID-only-trained variant of the model, not the submitted
generalist configuration).

\noindent\textbf{Ablation --- data mixing for generalization.} A FREUID-only
variant of the same architecture/recipe reached FREUID Score 0.00096 on the
public leaderboard but 0.9662 (AUC $\approx$0.51, i.e.\ chance level) on the
IDNet held-out set. Mixing in IDNet training data (this submission) reaches
0.1028 on the same held-out set --- an order-of-magnitude improvement in
measured cross-document-type generalization --- at a cost of public-leaderboard
score (0.00676), which we consider the correct trade given the private test
set's document-type composition.

\section{Reproducibility}
\label{sec:repro}
\begin{itemize}[noitemsep]
  \item \textbf{Repository:} \url{https://github.com/nadhirhasan/FREUID-Challenge-2026}, commit \texttt{b9473882f66a8baefce83980cf19f4db23d11a33}.
  \item \textbf{Docker:} \texttt{docker build -t freuid-repro:local .} from the repository root; weights are baked in at \texttt{/models/cv4\_fold0.pt} (Git LFS in the repository).
  \item \textbf{Dependencies:} Python 3.11, PyTorch 2.6.0 (+cu124), timm 1.0.27, opencv-python-headless 4.10.0.84 --- see \texttt{docker/requirements.txt} for the exact pinned inference environment (training used the same versions).
  \item \textbf{Runtime:} single RTX A4500 (20GB), $\approx$20 images/s at inference (batched, GPU); falls back to CPU automatically if no GPU is exposed to the container.
  \item \textbf{Randomness:} training uses unseeded per-worker RNGs for augmentation only (does not affect the frozen, deterministic inference path); inference is fully deterministic given fixed weights.
  \item \textbf{Network:} inference requires \texttt{--network none}; no downloads occur at container runtime. All dependencies and weights are installed/copied at build time.
\end{itemize}

\section{Team and contributions}
[FILL IN: team members and roles]

\section*{Acknowledgments}
We thank the FREUID Challenge organizers for the FantasyID/FREUID dataset and
the IDNet authors for their permissively licensed dataset.

\bibliographystyle{plain}
\bibliography{references}

\appendix
\section{Hyperparameters}
\label{app:hparams}
\begin{table}[h]
  \centering
  \small
  \begin{tabular}{ll}
    \toprule
    Hyperparameter & Value \\
    \midrule
    Image size & $448\times728$ (letterboxed) \\
    Batch size (train) & 6 (accum 4, effective 24) \\
    Batch size (eval) & 8 \\
    Learning rate (head) & $1\times10^{-3}$ \\
    Learning rate (LoRA) & $2\times10^{-4}$ \\
    Weight decay & 0.05 (head) / 0.0 (LoRA) \\
    LoRA rank / $\alpha$ / dropout & 16 / 32 / 0.05 \\
    Epochs & 3 \\
    Synthetic-attack rate (sbi) & 0.25 \\
    IDNet training cap & 56{,}000 (balanced) \\
    IDNet held-out (selection) size & 4{,}000 (balanced, 2 countries) \\
    Seed (fold split) & 42 \\
    \bottomrule
  \end{tabular}
\end{table}

\section{Additional ablations}
See Section~\ref{sec:results} for the two controlled single-lever ablations
(attack realism; FREUID-only vs.\ FREUID+IDNet data mixing) that motivated the
final configuration.

\end{document}
